% !TeX root = ../papers/nc-p3.tex
\section{半単純なBell直線の全補完・無印同型・最小周期}
\label{sec:split-spectral-line-classification}
\begin{theorem}\label{thm:split-spectral-line-classification}
定理\ref{thm:split-bell-locus}の直線の全点で、交代成分が半単純な可逆二平面分解窓を
添字保持同型で割ると、全ての量子多項式類と以下の四行で尽きる。全て標準Koszul左右AS正則で最小周期1である。
\begin{enumerate}
\item $s\in k^\times$, $s^2\ne1$ の一族、$s\sim-s\sim s^{-1}$。
\item $s^2=-1$ で追加される $a\in k^\times$ の一族、$a\sim-a\sim a^{-1}$。
\item $s^2=1$ で一方が非半単純Jordan、他方が半単純な一族。後者の固有値比
$\lambda\in k^\times$ を $\lambda\sim\lambda^{-1}$ で割り、$\lambda=1$ はスカラーとする。
\item $s^2=1$ で両方とも非半単純Jordanの一点。
\end{enumerate}
この四行は相互に交わらない。量子部分の全同型は定理\ref{thm:quantum-indexed-all}による。
\end{theorem}
\begin{proof}[印付きの全移送]
$P$ のBell四数は $(1,1,s,s)$ に取れ、付け替えによる $s$ の自由度は
$\pm s^{\pm1}$ である。通常の座標では $P_{abab}=K_{ab}$ が全ての非零成分で、
$K=\left(\begin{smallmatrix}1&s\\s&1\end{smallmatrix}\right)$ である。
両隣接切断の逆転置はどちらも同じ表示の $P(1/s)$ であり、
$\pi_0=(03)(12),\pi_1=(01)(23)$ の両方がこれを保つ。
第一の四対応を恒等に取れば、第二の四対応 $L_i$ の全体は
$P(1/s)$ の局所射影安定群 $H$ の元である。
$P(s)$ の安定群も $H$ と同じで、$b=(b_i)\in H$ による同型は
\[
                   L_i\longmapsto b_{i+2}^{-t}L_i b_i^t
\]
である。これは補空間側の基底変更を第一の四対応で消した式である。

$D(a)=\operatorname{diag}(1,a)$, $X=\left(\begin{smallmatrix}0&1\\1&0\end{smallmatrix}\right)$,
$Z=D(-1)$, $T(a,b)=(D(a),D(b),D(a^{-1}),D(b^{-1}))$ とする。
$s^2\ne1$ なら $02\mid13$ 切断の両支持は $00,11$ の張る二平面であり、
その中の可分解直線はこの二本だけである。ゆえに局所作用は、位置0,2で同じビット置換、
位置1,3でも同じビット置換を持つ単項行列である。係数表 $K$ の行列行・列倍率に関する不変量は
交差比 $s^{-2}$ である。片側だけの反転が可能なのは $s^4=1$ のときだけである。
$s^4\ne1$ なら $H$ は全 $T(a,b)$ と一斉反転 $F=(X,X,X,X)$ で生成される。
従って $L=T(a,b)$ は全て $L=1$ に同型、$L=T(a,b)F$ は
$(a,b)\sim(a^{-1},b^{-1})$ だけで割る。実際、前者には対角基底の二乗が作用し、
後者で対角基底は自明に作用する。

$s^2=-1$ では二つの独立な反転の持上げを
\[
           A=(XZ,Z,X,1),\qquad B=(Z,XZ,1,X)
\]
に取る。これらは局所射影群で可換な対合であり、$AB=F$ である。
全 $L$ は $T(a,b)A^e B^f$ ($e,f\in\{0,1\}$) となる。
対角基底で第一座標は $e=0$ のとき、第二は $f=0$ のとき自由に消せる。
残る座標で $A,B$ による作用は、それぞれ
$(a,b)\mapsto(-a^{-1},-b)$ と $(a,b)\mapsto(-a,-b^{-1})$ である。
ここでは消した座標を無視し、$e=f=0$ なら一点となる。
従って片側だけの反転は各側で一つの $a\sim-a\sim a^{-1}$ の族となる。
両側反転 $e=f=1$ では表示した二つの作用で二座標を割る。

$s=1$ に残る場合は $K$ の階数が1で、$P$ は $02,13$ の二つの非退化形式の積である。
このとき $H$ は $(C,D,C^{-t},D^{-t})$ の全体で $\mathrm{PGL}_2^2$ と同型である。
$L$ の最初の二行列 $C,D$ を各々独立な射影合同で割ればよい。
$J=\left(\begin{smallmatrix}0&-1\\1&0\end{smallmatrix}\right)$ とすると
これは $CJ,DJ$ の独立な射影共役に等しい。
\end{proof}
\begin{proof}[量子窓の認識と印の除去]
両側反転 $L=T(a,b)F$ なら各位置2,3の全二文字でビットを反転すると、
四テンソルの24個の非零語はすべて四文字の置換となる。
逆Schur式の倍率は全て非零である。各辺の切断は、互いに補となる無順序二文字組について
六つの二行二列ブロックに分かれ、全階数6、各係数非零だから各ブロックが階数1である。
よって全関係はスカラー付きの二文字交換である。三次PBW次元から相異なる三文字を
逆順から並べる両経路の係数が一致し、各三角積は位置に依存しない。位置別の対角倍率で
交換係数を一定にできる。具体的に $L=T(a,b)F$ での四三角積は、文字を二色の順に並べて
$(-as,-a/s,-1/(bs),-s/b)$ である。\path{computations/verify_split_spectral_line.py} は両切断の再構成から
全位置でこの式を直接求め、例外点の基底変換と一移動も検算する。$s=1$ で $CJ,DJ$ がともに半単純な場合も、各々を対角化すれば
$C,D$ が反対角になり同じ議論が通る。逆に任意の量子窓で四座標を二組に分けると
射影座標二直線が指定できる。その交代四テンソルは位置別のビット変更の後 $abab$ にのみ支持され、
四係数は全て非零である。行・列倍率で対角が1、非対角が共通の $s\ne0$ にできるので
直線スペクトルに入る。

分解を与える一径数群を含む極大トーラスの次元は、共通法線倍率の1に、
$H$ 内の $L$ の上記移送作用による安定群の階数を足したものである。
実際この一径数群の中心化群は分解を保ち、第一の四同定を固定すれば
残る相対倍率の核は六つの二重次数から共通の法線倍率だけになる。
$s^2\ne1$ の反転のない場合は階数1、片側だけの反転なら2である。
$s=1$ で非半単純な行列の射影中心化群の恒等成分は加法群なので、一つだけ非半単純なら階数2、
両方なら1である。量子窓の階数は少なくとも3なのでこれらは量子部分と交わらない。

非量子の場合の二平面分解を調べる。階数1なら分解を与える原始余指標は共役と符号を除いて一つである。
階数2では共通法線倍率のほか、片側の二文字だけを非スカラーに動かすトーラスがある。
これは $e=1,f=0$ なら $T(a,1)$、$s=1$ で一方だけ非半単純なら
半単純側の対角トーラスである。極大トーラス内の余指標が二文字ずつの二重みを持つには、
スカラーで作用される側の二文字が既に同じ重みなので、残る二文字も同じ重みでなければならない。
よってやはり共通法線倍率またはその逆である。
従って無印同型で追加されるのは分解の交換だけである。

交換では第一の四対応は恒等、第二は $L_{\pi_1(i)}^t$ となる。
交代成分 $P(1/s)$ を $P(s)$ に戻すには新しい能動位置0,2で $X$ を作用させればよい。
これで $e,f$ が入れ替わり、反転なしの一点は同じ一点、片側反転の二族は互いに移る。
片側反転を位置0,2に選んで比較すれば、追加の自己比較はなく第二行の同値が残る。
$s=1$ では交換に際して $P$ の変更は不要で、射影合同類の対が入れ替わる。
二次行列とその転置は射影合同である。実際 $C^t J$ は $CJ$ の逆行列と
射影共役になり、二次元射影行列はその逆と共役である。
従って最後の二行を得る。両側の取替えと逆行列は $02\mid13$ の支持階数も保つので
$s^2=1$ と他の非量子行は交わらない。
\end{proof}
\begin{proof}[一移動]
一移動と同時に分解を交換すると、新しい能動四テンソルは $P$ の一巡回で、
$K$ が対称だから元の $P$ に等しい。第一の四対応は旧第二の $L_{i+1}$ であり、
これを恒等に戻した後の第二の四対応は $L_{i-1}^{-1}$ である。
反転なしでは $L=1$ が戻る。片側反転 $L=T(a,1)A$ ではここから再び分解を交換し、
能動位置0,2の $X$ で $P$ を戻すと、第二の四対応は
$X^{[i\ {\rm even}]}L_{\pi_1(i)-1}^{-t}X^{[i\ {\rm even}]}$ である。
偶位置は元と同じで、奇位置の $Z,1$ が入れ替わるだけであり、自由に消せる対角係数で元に戻る。
$s=1$ の一移動と交換では、先頭の二行列が $D^t,C^{-1}$ になる。
$C^{-1}$ は $C^t$ と射影合同なので無印の対は等しい。量子部は一定の次数付き表示を持つ。
これで全点の最小周期が1である。
\end{proof}
